\chapter*{Preface}
The second part of \emph{Order and Chance} builds machines: measure and integration, Hilbert spaces and operators, the
Fourier transform on groups, Haar measure, manifolds and homology, Lie groups, commutative and homological algebra. A machine
is learned when it has been taken apart and put together again, and that is what the four volumes of Part~II do. This book is
about switching the machines on. Its three hundred problems take questions that could be asked in Part~I, about fractions,
primes, polynomials, random walks and permutations, and answer them with the machinery of Part~II, or take a machine and ask
what it says about numbers.

The questions are often old and simple. How well can a typical real number be approximated by fractions? How many ways are
there to cut a necklace fairly between two thieves? Why must a closed curve in space turn through at least a full circle,
and twice that if it is knotted? How many ways are there to write a cycle as a product of transpositions? The answers use
the Borel--Cantelli lemma and Hausdorff dimension, the Borsuk--Ulam theorem, integral geometry on the sphere, and the
characters of the symmetric group. Chapters are organised by meeting points, never by volumes, and the statements never say
which tool applies: finding it is half of each problem.

The format is that of the first problem book. Problems are graded \textbf{$\star\star$} for a good afternoon and
\textbf{$\star\star\star$} for a gem that may take a weekend; a few \textbf{$\star$} problems open each section, and the
problems marked \emph{Journey} are long walks in many steps to a famous theorem. Hints, kept apart at the end of each chapter,
give a nudge when the search stalls. Every solution starts with \emph{How to find it}, which tells how one might think of the
approach, continues with a complete \emph{Write-up}, names the \emph{Tools used} with their volumes and chapters, and often
ends with \emph{Where it leads}, where research results quoted for their interest are labelled as such and never used in a
proof. Everything that is used is proved here or in a named volume.


The pleasure of Part~II is that abstraction pays. A theorem about all Hilbert spaces, or all compact groups, turns out to know
something precise about the integers. We have tried to collect as many of those surprises as possible.
